\documentclass[aspectratio=169]{beamer}
\usetheme{Bruno}
\usepackage{amsmath}
\usepackage{siunitx}
\usepackage[american,RPvoltages]{circuitikz}
\ctikzset{capacitors/scale=0.7}
\usepackage{tabularx}
\newcolumntype{C}{>{\centering\arraybackslash}X}
\usepackage{tabu}
\usepackage[spanish]{babel}
\usepackage{booktabs}
\usepackage{pgfplots}
\title{Electricidad I \\ \emph{Circuitos de primer orden}}

\usetikzlibrary{arrows, arrows.meta}
\author{
    Juan J. Rojas
}

\institute{Instituto Tecnológico de Costa Rica}
\date{\today}
\background{fig/background.jpg}
\renewcommand\tabularxcolumn[1]{m{#1}}% for vertical centering text in X column


\begin{document}

\maketitle
\sisetup{unit-math-rm=\mathrm,math-rm=\mathrm} % change sinitx font

\begin{frame}{Circuito RC: respuesta natural}
    \begin{columns}
    \begin{column}{0.5\textwidth}
        \centering
        \only<1>{
            \begin{circuitikz}[scale=0.8]\draw
            (0,2.5) node[spdt, rotate = 90, anchor=in] (Sw) {}
            (Sw.in) 
                to [C, l=$C$, voltage shift=3]
            (0,1) -- (0,0)
            (Sw.out 1) |- (-1,4) to [R,l=$R_f$,o-] (-4,4)
                to[V, l_=$V$, invert]
            (-4,0)
            (Sw.out 2) |- (1.5,4) to[short,o-] (3,4)
                to[R,l=$R$]
            (3,0) to[short,-o] (-1,0) to[short,-o] (1.5,0) -- (-4,0)
            (-1,4)
                to[open,v=$v$,voltage shift=-3]
            (-1,0)
            ;
        \end{circuitikz}
        }
        \only<2>{
            \begin{circuitikz}[scale=0.8]\draw
            (0,2.5) node[spdt, rotate = 90, yscale = -1, anchor=in] (Sw) {}
            (Sw.in) 
                to [C, l=$C$, voltage shift=3]
            (0,1) -- (0,0)
            (Sw.out 2) |- (-1,4) to [R,l=$R_f$,o-] (-4,4)
                to[V, l_=$V$, invert]
            (-4,0)
            (Sw.out 1) |- (1.5,4) to[short,o-] (3,4)
                to[R,l=$R$]
            (3,0) to[short,-o] (-1,0) to[short,-o] (1.5,0) -- (-4,0)
            (1.5,4)
                to[open,v=$v$,voltage shift=-3]
            (1.5,0)
            ;
        \end{circuitikz}
        }
    \end{column}
    \begin{column}{0.5\textwidth}
        \only<1>{
        Después de mucho tiempo...
        \begin{equation*}
            v = V
        \end{equation*}
        }
        \only<2>{
        Si se cambia la posición del interruptor en $t=0$
        \begin{equation*}
            v(0) = V, t_0 = 0
        \end{equation*}
        }
    
    \end{column}
\end{columns}
\end{frame}

\begin{frame}{Circuito RC: respuesta natural}
    \begin{columns}
    \begin{column}{0.5\textwidth}
        \centering
        \only<1-12>{
            \begin{circuitikz}[scale=0.8]\draw
            (0,4)
                to [C, l=$C$, i>^=$i_C$]
            (0,0)
            (0,4) -- (2,4) to[short,o-] (4,4)
                to[R,l=$R$, i>^=$i_R$]
            (4,0) to[short,-o] (2,0) -- (0,0)
            (2,4)
                to[open,v=$v$,voltage shift=-3]
            (2,0)
            ;
        \end{circuitikz}
        }
        \only<13-15>{
            \begin{tikzpicture}
            \begin{axis}[
                %yshift=-1cm,
                width= \linewidth,
                height= 0.7\linewidth,
                %grid=major, % Display a grid
                %grid style={dashed,gray!30}, % Set the style
                %every axis plot/.append style={thick},
                ymin=0, ymax=1.1,
                xmin=0, xmax=6,
                axis x line=bottom,
                axis y line=left,
                axis line style={-},
                xtick=\empty,
                xticklabels=\empty,
                ytick=\empty,
                yticklabels=\empty,
                xlabel=Voltage, % Set the labels
                ylabel=Current,
                x unit=V, 
                y unit=A,
                %legend pos=north east,
                %legend style={cells={anchor=west}}
                %legend style={at={(1,0.7)},anchor=west}, % Put the legend below the plot
                smooth,
                ]
                \addplot [name path=VI, mark=none, color=black,forget plot, thick]       
                table[x=t,y=mag,col sep=comma] {./python/resp_nat.csv};
            \end{axis}
            \end{tikzpicture}
        }
    \end{column}
    \begin{column}{0.5\textwidth}
    \only<1-3>{
        \onslide<1-3>{
        Aplicando LCK tenemos que:
        \begin{equation*}
            i_C + i_R = 0
        \end{equation*}
        }
        \onslide<2-3>{
        y sabemos que: 
        \begin{equation*}
            i_C = C\,\frac{dv}{dt},\, i_R=\frac{v}{R}
        \end{equation*}
        }
        \onslide<3>{
        entonces:
        \begin{equation*}
            C\,\frac{dv}{dt} + \frac{v}{R} = 0
        \end{equation*}
        }
    }
    \only<4-6>{
        \onslide<4-6>{
        dividiendo entre $C\,$ y reordenando:
        \begin{equation*}
            \frac{dv}{dt} = -\frac{v}{RC}
        \end{equation*}
        }
        \onslide<5-6>{
        reordenando nuevamente:
        \begin{equation*}
            \frac{1}{v} dv = -\frac{1}{RC} dt
        \end{equation*}
        }
        \onslide<6>{
        integrando a ambos lados:
        \begin{equation*}
            \int\limits_{v(0)}^{v(t)}\frac{1}{v} dv = \int\limits_{t_0}^{t}-\frac{1}{RC} dt
        \end{equation*}
        }
    }
    \only<7-9>{
        \onslide<7-9>{
        tenemos que:
        \begin{equation*}
            \ln v\,\bigg|_{v(0)}^{v(t)} = -\frac{t}{RC}\,\bigg|_{t_0}^{t}
        \end{equation*}
        }
        \onslide<8-9>{
        evaluando:
        \begin{equation*}
            \ln v(t) - \ln v(0) = -\frac{1}{RC}\left(t - t_0\right)
        \end{equation*}
        }
        \onslide<9>{
        entonces:
        \begin{equation*}
            \ln\frac{v(t)}{v(0)} = -\frac{1}{RC}\left(t - t_0\right)
        \end{equation*}
        }
    }
    \only<10-12>{
        \onslide<10-12>{
        eliminamos el logaritmo natural:
        \begin{equation*}
            e\,^{\ln \left( v(t)/v(0)\right)} = e\,^{-(t-t_0)/RC}
        \end{equation*}
        }
        \onslide<11-12>{
        y obtenemos:
        \begin{equation*}
            v(t) = v(0) e\,^{-(t-t_0)/RC}
        \end{equation*}
        }
        \onslide<12>{
        pero sabemos que $t_0 = 0$\, y $v(0)=V$ entonces:
        \begin{equation*}
            v(t) = Ve\,^{-t/RC}
        \end{equation*}
        }
    }
    \only<13-15>{
        \onslide<13-15>{
       la velocidad a la que la tensión decrese $RC$:
        \begin{equation*}
            e\,^{\ln \left( v(t)/v(0)\right)} = e\,^{-(t-t_0)/RC}
        \end{equation*}
        }
        \onslide<14-15>{
        y obtenemos:
        \begin{equation*}
            v(t) = v(0) e\,^{-(t-t_0)/RC}
        \end{equation*}
        }
        \onslide<15>{
        pero sabemos que $t_0 = 0$\, y $v(0)=V$ entonces:
        \begin{equation*}
            v(t) = Ve\,^{-t/RC}
        \end{equation*}
        }
    }
    \end{column}
\end{columns}
\end{frame}



\begin{frame}{Inductores en serie}
\begin{columns}
    \begin{column}{0.5\textwidth}
    \begin{center}
        \begin{circuitikz} [scale=1]\draw
            (3,0)
                to[L,l=$L_3$,o-]
            (0,0)	
                to[L,l=$L_2$]
            (0,3)
                to[L,l=$L_1$,i^<=$i$,-o]
            (3,3)
            ;
        \end{circuitikz}
    \end{center}
    \end{column}
    \begin{column}{0.5\textwidth}
    \begin{center}
        \begin{equation*}
            L_{eq}=L_1+L_2+L_3
        \end{equation*}
    \end{center}
    \end{column}
\end{columns}
\end{frame}

\begin{frame}{Inductores en paralelo}
\begin{columns}
    \begin{column}{0.5\textwidth}
    \begin{center}
        \begin{circuitikz} [scale=1]\draw
            (1,3)
                to[short,o-]
            (0,3)	
                to[L,l_=$L_3$]
            (0,0)
                to[short,-o]
            (1,0)
            (0,3) -- (-1.5,3)
                to[L,l_=$L_2$]
            (-1.5,0) -- (0,0)
            (-1.5,3) -- (-3,3)
                to[L,l_=$L_1$]
            (-3,0) -- (-1.5,0)
            (1,3)
                to[open,v^=$v$]
            (1,0)
            ;
        \end{circuitikz}
    \end{center}
    \end{column}
    \begin{column}{0.5\textwidth}
    \begin{center}
        Caso general:
        \begin{equation*}
            \frac{1}{L_{eq}}=\frac{1}{L_1}+\frac{1}{L_2}+\frac{1}{L_3}
        \end{equation*}
       Caso especial, solo dos inductores:
        \begin{equation*}
            L_{eq}=\frac{L_1\cdot L_2}{L_1+ L_2}
        \end{equation*}
    \end{center}
    \end{column}
\end{columns}
\end{frame}

\begin{frame}{Resumen}
    \begin{center}
        \begin{tabularx}{14cm}{C C C}
        \toprule
        Variable & Ecuación $v-i$ & Ecuación $i-v$ \\
        \midrule
       Capacitor & $v=\frac{1}{C}\int_{t_0}^{t} i(t)dt + v(t_0) $ & $i = C\,\frac{dv}{dt}$ \\[5pt]
        Inductor & $v = L\,\frac{di}{dt}$ & $i=\frac{1}{L}\int_{t_0}^{t} v(t)dt + i(t_0) $\\[5pt]
        \bottomrule
        \end{tabularx}   
    \end{center}
\end{frame}

% \begin{frame}{Referencias}

% \bibliographystyle{ieeetr}

% \bibliography{referencias}

% \end{frame}

\end{document}